%%%%%%%%%%%%%%%%%%%%%%%%%%%%%%%%%%%%%%%%%%%%%%%%%%%%%%%%%%%%%
%%%%%% Article draft for submission to ISeCure Journal %%%%%%
%%%%%%%%%%%%%%%%%%%%%%%%%%%%%%%%%%%%%%%%%%%%%%%%%%%%%%%%%%%%%
\documentclass[Print]{../Style/isecure-v24}

\pdfminorversion=4
%\usepackage{amsmath,amssymb}
%\usepackage{array}
\usepackage{lipsum}
\usepackage{algpseudocode,algorithm}
\usepackage[colorlinks, bookmarksnumbered, linkcolor=darkblue, citecolor=darkred, urlcolor=darkgreen]{hyperref}
\usepackage[numbers,sort&compress]{natbib}
\usepackage[all]{hypcap}
\usepackage{../Style/picins}
\usepackage{lettrine}
\usepackage{wrapfig}
\usepackage{graphicx}
\usepackage{subcaption}
%\usepackage{multicol}

%\hypersetup{colorlinks=true}
\usepackage{amsthm}
\usepackage{amssymb,amsmath,amsfonts,eqnarray,breqn}
\usepackage{fancyhdr}
%\usepackage{cite}
\usepackage{ragged2e}
\usepackage{enumitem}   
\usepackage{refcount}
\usepackage[protrusion=true,expansion=true]{microtype}
\usepackage{float}
\usepackage[greek,english]{babel}
\usepackage[LGR,T1]{fontenc}
%%%%%%%%%%%%%%%%%%%%%%%%%%%%%%%%%%%%%%%%%%%%%%%%%%%%%%%%%%%%%%%%%%%%%%%%%
%\theoremstyle{plain}

\DisableLigatures[f]{encoding=*,family=*}
\graphicspath{{Images/},{../Style/}}
\DeclareGraphicsExtensions{.jpg,.pdf,.png}
\input{../Style/Theorem-Styles.tex}
\articletype{Short Paper}
\def\sectionautorefname{Section} % localizing \autoref{}
\def\subsectionautorefname{Section} % localizing \autoref{}
\def\subsubsectionautorefname{Section} % localizing \autoref{}
\providecommand*{\lemmaautorefname}{Lemma}
%\newcommand{\subsubsectionautorefname}{Section}
\usepackage{titlesec}
\setcounter{secnumdepth}{4}
\titleformat{\paragraph}
{\normalfont\normalsize\bfseries}{\theparagraph}{1em}{}
\titlespacing*{\paragraph}
{0pt}{3.25ex plus 1ex minus .2ex}{1.5ex plus .2ex}

\journal{ISeCure}
\company{ISC}
\copyrightyear{2018}
\journalsite{http://www.isecure-journal.org}
\firstpage{1} % TODO



\makeatletter
\begingroup \lccode`+=32 \lowercase
 {\endgroup \def\Url@ObeySp{\Url@Edit\Url@String{ }{+}}}
 \def\Url@space{\penalty\Url@sppen\ }
\makeatother


%%%%% Use the following commands for putting a watermark
%\usepackage{type1cm}
%\usepackage{eso-pic}
%\makeatletter
%\AddToShipoutPicture{
%            \setlength{\@tempdimb}{.5\paperwidth}
%            \setlength{\@tempdimc}{.5\paperheight}
%            \setlength{\unitlength}{1pt}
%            \put(\strip@pt\@tempdimb,\strip@pt\@tempdimc){
%        \makebox(600,-900){\rotatebox{45}{\textcolor[gray]{0.90}
%        {\fontsize{2cm}{2cm}\selectfont{\bfseries{Uncorrected~~~~Proof}}}}}
%            }
%}
%\makeatother

%Use the following command for putting line numbers
\usepackage[displaymath]{lineno}  %DON'T COMMENT EVEN IF YOU DONT WANT LINE NOS
\newcommand*\patchAmsMathEnvironmentForLineno[1]{
  \expandafter\let\csname old#1\expandafter\endcsname\csname #1\endcsname
  \expandafter\let\csname oldend#1\expandafter\endcsname\csname end#1\endcsname
   \renewenvironment{#1}
     {\linenomath\csname old#1\endcsname}
     {\csname oldend#1\endcsname\endlinenomath}}
\newcommand*\patchBothAmsMathEnvironmentsForLineno[1]{
  \patchAmsMathEnvironmentForLineno{#1}
  \patchAmsMathEnvironmentForLineno{#1*}}
\AtBeginDocument{
\patchBothAmsMathEnvironmentsForLineno{equation}
\patchBothAmsMathEnvironmentsForLineno{align}
\patchBothAmsMathEnvironmentsForLineno{flalign}
\patchBothAmsMathEnvironmentsForLineno{alignat}
\patchBothAmsMathEnvironmentsForLineno{gather}
\patchBothAmsMathEnvironmentsForLineno{multline}
}

%%%%%%%%%%%%%%%%%%%%%%%%%%%%%%%%%%%%%%%%%%%%%%%%%%%%%%%%%%%%%%%%%%%%%%%%

%%%%%%%%%%%%%%%%%%%%%%%%%%%%%%%%%%%%%%%%%%%%%%%%%%%%%%%%%%%%%%%%%%%%%
\newtheorem{cor}{Corollary}[section]
\newtheorem{exmp}{Example}[section]
\newtheorem{construction}{Construction}[section]
\newtheorem{defn}{Definition}[section]
\newcommand*\concat{\mathbin{\|}}
%%%%%%%%%%%%%%%%%%%%%%%%%%%%%%%%%%%%%%%%%%%

\begin{document}
%\linenumbers
\begin{frontmatter}

%------------------- Paper Title ------------------------------
% Titles are generally capitalized except for words such as a, an, and, as, at, but, by, for, in, nor, of, on, or, the, to and up, which are usually not capitalized unless they are the first or last word of the title. Linebreaks \\ can be used within to get better formatting as desired. Do not put math or special symbols in the title. Title of up to 17 words should not contain the name of locations, countries or cities of the research as well as abbreviations. Avoid complicated and technical expressions and do not use vague expressions. The title should be oriented to ISeCure scope issues while not being obscure or meaningless.

\def\NoDingTitle{Manuscript Title}
\title{\NoDingTitle\textsuperscript{ }}

%------------------- Author names and Affiliations ------------------------------
% Add your name, affiliations and email in this part. Please consider that an academic email is preferred. If the fisrt author isn't corresponding author in your case, please set the CorAuth option in \author command for that author. 

\author[a1,CorAuth]{Author 1}, and%
\ead{ Author1@email.com}
\author[a2]{Author 2}
\ead{Author1@email.com}


\address[a1]{Author1 affiliation}
\address[a2]{Author2 affiliation}

\corauth[CorAuth]{Corresponding author.}% Do not change this field.



%------------------- Abstract ------------------------------
\begin{abstract}
An abstract of 150 to 250 words that sketches the purpose of the study, basic procedures, main findings, its novelty, discussions and the principal conclusions goes here. Moreover, The abstract should not contain any math formula, undefined abbreviations, and references.
%% The following command generate a random text.
\lipsum[2]
\end{abstract}

%---------------------- Keywords -------------------------------
% Please provide 3 to 5 keywords, each separated by hyphen and started in capital letter.
\begin{keyword}
Provide 4 to 7 keywords which can be used for indexing purposes. Keywords should not repeat the words of the manuscript title or contain abbreviations and shall be written in alphabetical order as separated by Comma and started in capital letter. Abbreviations should be defined at first mention and used consistently thereafter through the text.

\end{keyword}


\makeatother

\end{frontmatter}

\addtolength{\parskip}{2mm}

%---------------------- Introduction ---------------------------

\section{Introduction} \label{sec:intro}

First section of the manuscript should be an introduction of the proposed work. Others vary according to the subject and contributions of the author. Please be noted that you should have Introduction, Related Work and Conclusion sections in your manuscript. In the introduction, briefly discuss the existing problem and your contribution in 3 to 5 paragraphs. The last paragraph should be an overview on the following sections and their contents. 

\lipsum[2-5]

\section{Related Work} \label{sec:related}
In this section, you should summarize and classify the related work. 

\lipsum[6-8]

\section{Rules} \label{sec:rules}
We appreciate that you have selected our journal for publishing your paper. We set up some rules to standardize the papers we receive in order to make it easier for the author and the editor. Please consider these rules in writing your manuscript.
\subsection{Images} \label{sec:images}

Figures/ illustrations should be in high quality art work, within 200-300 dpi and separately provided. Ensure that figures are clear, labeled, and of a size that can be reproduced legibly in the journal. Figures should be set within the text. Following remarks should be applied to the figures:

\begin{itemize}

\item Each figure should have a concise caption describing accurately what the figure depicts. Include the captions in the text file of the manuscript, not in the figure file.

\item Figure captions begin with the term Fig. Figures should be with the captions placed below in limited numbers.

\item No punctuation is to be placed at the end of the caption. 

\item Identify all elements found in the figure in the figure caption; and use boxes, circles, etc., as coordinate points in graphs.

\item Identify previously published material by giving the original source in the form of a reference citation at the end of the figure caption.

\end{itemize}

Use the following format for your figures and cite to each with command \textbackslash autoref$\{figure\_label\}$ such as~\autoref{fig:sample}. Please copy your image into the Image folder.
\begin{figure}
	\centering
	\includegraphics[width=0.2\textwidth]{sample.png} % image file name
	\caption{Sample Figure}
	\label{fig:sample} % your defined label.
\end{figure}

\subsection{Tables} \label{sec:tables}

Do not submit tables and graphs as photograph. Tables should be set within the text. Do not use internal horizontal and vertical rules. Tables should be called out in the text and should have a clear and rational structure and consecutive numerical order. All tables should be numbered (1, 2, 3, etc.). Give enough information in subtitles so that each table is understandable without reference to the text. For each table, please supply a table caption (title) explaining the components of the table. Identify any previously published material by giving the original source in the form of a reference at the end of the table caption. Tables should be with the captions placed above in limited numbers.

Footnotes to tables should be indicated by superscript lower-case letters (or asterisks for significance values and other statistical data) and included beneath the table body.

Use the following format for tables and cite to each with command \textbackslash autoref$\{table\_label\}$ such as ~\autoref{tab1}.

\begin{table}
	\tiny
	\centering
	\caption{A sample table}\label{tab1}
	\begin{tabular}{ccccc}
		\cline{2-5} 
		1~~~~~ & 2~~~~ & 3~~~~   & 4~~~~~ & 5 ~~~~   \\ 
		\hline 
		$E$ & $S^{ID}$ &  $-$ &  $-$ &  $-$    \\  
		$P_1$ & $IS^1$ & $-$  & $-$   &  $-$   \\ 
		$P_2$ & $IS^2$ &  $-$ & $-$  & $-$  \\ 
		$\vdots$ & $\vdots$ & $\vdots$  & $\vdots$  & $\vdots$   \\  
		$P_n$ & $IS^n$ & $-$  &  $-$ &  $-$  \\ 
		\hline 
	\end{tabular} 
\end{table}


\subsection{Algorithms and Equations} \label{sec:algos}
	
We suggest you to use the ``algpseudocode'' and ``algorithm'' packages to add your algorithms and pseudo codes to the manuscript. You can find an example in Pseudocode~\ref{alg}.

\begin{algorithm}
	\floatname{algorithm}{Pseudocode}
	\caption{General Behavioral Patterns Detection Rules}
	\label{alg}
	%\begin{algorithm}
	\begin{algorithmic}[1]
		\Function{DetectTryToConnect}{$profiles$}
		\State $TC \gets$ all profiles that $SendCount=0$ and  $ReceiveCount=0$
		\State $D[i] \gets$ number of profiles in $TC$ where  $Distance=i$
		\State $M \gets$ maximum value in array $D$
		\If {$M \ge \delta_1$}
		\Return \textsf{TRY\_TO\_CONNECT}
		\EndIf
		\EndFunction		
	\end{algorithmic}
\end{algorithm}
	
It is recommended to use the following format for numbering the equations:

\begin{equation}
	a= b\oplus c
\end{equation}

To remove the number use the following format:

\begin{equation*}
	a= b\oplus c
\end{equation*}

For a set of Equations use the following format:

\begin{align*}
	a &= b\oplus c\\
	a &= b\oplus c\\
	a &= b\oplus c\\
	a &= b\oplus c
\end{align*}

\subsection{Theorems, Lemma, Proofs} \label{sec:theorems}

Use the following format for theorems, lemma, and proofs and cite to each with command \textbackslash autoref$\{theorem\_label\}$ such as~\autoref{theorem1}.


\begin{theorem} \label{theorem1}
	Here is the theorem description.
\end{theorem}
\begin{proof}
	Write the theorem proof here.
\end{proof}

\begin{lemma}
	Here is the lemma description.
\end{lemma}

\begin{proof}
	Write the lemma proof here.
\end{proof}

\subsection{Sections and Subsections}

Use the following command to create a new sections and subsections:

\begin{itemize}
	\item Section: ``\textbackslash section $\{Section\_Name\}$'
	\item Subsection: ``\textbackslash subsection $\{Subsection\_Name\}$'
\end{itemize}

And easily refer to each of them using ``\textbackslash ref$\{Label\}$'' with the label you have defined using ``\textbackslash label$\{label\}$''. For example to reference to Introduction Section~\ref{sec:intro}.

\subsection{Citation and References}

To cite a paper please use ``\textbackslash cite$\{\}$'' such as~\cite{grumbach2017distributed}. For more than one citation simply separate them with comma in  ``\textbackslash cite$\{\}$''~\cite{lenstra2015random, he2009electronic}. However, avoid citing to many papers at once. The tilde character (\~{}) in the source means a non-breaking space. This way, your reference will always be attached to the word that preceded it, instead of going to the next line.

We highly suggest you to add your references in a ``.bib'' file.  The default bib file is ``biblio.bib''. In case you use a file with another name, please include it with command~\\bibliography$\{\}$ in the main file.

If you intend to use ``bibliography'', please be noted that the ``MLA'' format  of the \href{https://scholar.google.com/schhp?hl=en&as_sdt=0,5}{Google scholar} is only accepted.


\section{Conclusion}\label{sec:conclusion}
This section should highlight the major, firm discoveries, and state what the added value of the main finding is, without literature references.

\section*{Acknowledgment}

Acknowledgments of people, grants, funds, etc. should be placed in a separate section before the reference list. The names of funding organizations should be written in full. Financial support affiliation of the study, if exists, must be mentioned in this section. Thereby, the Grant number of financial support must be included.

%%%%%%%%%%%%%%%%%%%%%%%%%%%%%%%%%%%%%%%%%%%%%%%%%
%%%%%%%%%%%%%%%%%%%%%%%%%%%%%%%%%%%%%%%%%%%%%%%%%
\bibliographystyle{unsrt}
\bibliography{biblio}
%%%%%%%%%%%%%%%%%%%%%%%%%%%%%%%%%%%%%%%%%%%%%%%%%
%%%%%%%%%%%%%%%%%%%%%%%%%%%%%%%%%%%%%%%%%%%%%%%%%


\begin{wrapfigure}[8]{l}{1.5cm}
	\begin{center}
		\includegraphics[width=2.3cm, height= 2.85cm]{Images/1.jpg}
	\end{center}
\end{wrapfigure}
\noindent
\newline
{\bf Author first name last name} {A short biography about the author. \lipsum[10]}

\begin{wrapfigure}[8]{l}{1.5cm}
	\begin{center}
		\includegraphics[width=2.3cm, height= 2.85cm]{Images/1.jpg}
	\end{center}
\end{wrapfigure}
\noindent
\newline
{\bf  Author first name last name} {A short biography about the author. \lipsum[11]}


\end{document}
